\begin{abstract}
\ifdefined\toyimageonly
We extend continuous bitstream diffusion (CoBit) from language modeling to multimodal
generation through a common binary interface. Frozen tokenizers or direct binary encodings map
observations into a joint bitstream, which CoBit learns to generate by reversing Gaussian
corruption of continuous bit representations. For each modality combination, we train one
denoiser from scratch, with a shared input projection, transformer backbone, bitwise prediction
head, and denoising objective. The same model generates the modalities jointly or generates
missing modalities conditioned on observed ones, which remain fixed during denoising. CoBit's
entropy-derived schedule adapts training noise levels and sampling steps to the data's estimated
information-loss profile. We first validate the approach on MNIST-Sum, a controlled image--text
task pairing raw binary images with equations encoded as text, achieving near-perfect conditional
accuracy and high joint consistency without learned tokenizers. We then test the same recipe at
larger scale on speech--text and protein sequence--structure modeling. Trained without text
pretraining, a single 633M-parameter speech--text model performs speech synthesis, recognition,
continuation, and unconditional co-generation. Jointly generated speech and text achieve 3.5\%
word error rate between the generated text and automatic transcriptions of the generated speech,
alongside a UTMOS speech naturalness score of 3.77. In protein modeling, one denoiser performs
folding, inverse folding, and joint sequence--structure generation. Generated backbones spanning
100--500 residues achieve a mean self-consistency TM-score of 0.899 after sequence redesign and
refolding, with 68.8\% scoring at least 0.9, demonstrating computational backbone designability.
These results establish the applicability of a common binary interface and bitstream diffusion to
multimodal generation.
\else
\note{Luca will lead the pitch and framing of the paper; to be written last, once results are final.}
\fi
\end{abstract}
